\documentclass[11pt]{article}

\usepackage[margin=1in]{geometry}
\usepackage{amsmath,amssymb,amsthm}
\usepackage{mathtools}
\usepackage{enumitem}
\usepackage{booktabs}
\usepackage{array}
\usepackage{graphicx}
\usepackage{tikz}
\usetikzlibrary{arrows.meta, positioning, calc}
\usepackage{hyperref}
\usepackage{cleveref}
\usepackage{algorithm}
\usepackage{algpseudocode}
\usepackage{xcolor}

\hypersetup{
    colorlinks=true,
    linkcolor=blue!50!black,
    citecolor=blue!50!black,
    urlcolor=blue!50!black,
}

\newtheorem{theorem}{Theorem}
\newtheorem{proposition}[theorem]{Proposition}
\newtheorem{lemma}[theorem]{Lemma}
\newtheorem{corollary}[theorem]{Corollary}
\newtheorem{definition}[theorem]{Definition}
\newtheorem{remark}[theorem]{Remark}
\newtheorem{example}[theorem]{Example}


\title{A Closed-Form, Non-Recursive Move Formula for the\\
Tower of Hanoi, and Its Behavior on Non-Canonical\\
Initial Configurations}

\author{}
\date{}

\begin{document}
\maketitle

\begin{abstract}
The Tower of Hanoi is almost universally solved by a two-line recursive
algorithm. We revisit a closed-form alternative: a formula that, given only
a move index $m$, the disk count $n$, and a source--destination peg pair
$(s,f)$, returns the disk moved on turn $m$ and the peg pair it travels
between --- with no recursion, no stack, and no simulated state. We give a
self-contained derivation of this formula from first principles, showing it
is equivalent to two classical facts about the optimal Hanoi sequence: the
\emph{ruler sequence} that governs which disk moves on turn $m$
(OEIS A001511), and the fact that the smallest disk cycles around the three
pegs in a single fixed rotational direction determined by the parity of $n$
and by $(s,f)$. We then examine an extension used by the reference
implementation to handle \emph{arbitrary} legal starting configurations
(disks distributed arbitrarily over the three pegs, not merely stacked on
one), by applying the closed-form formula disk-by-disk from the inside out.
We prove this construction always terminates in a valid sequence of legal
moves, but we show empirically that it is frequently and sometimes severely
suboptimal on scattered configurations: in a sample of $300$ random legal
$3$--$5$-disk configurations, roughly one third were solved with strictly
more moves than the true graph-distance optimum, with one case using
$51$ moves against an optimal value of $24$. We give an exact formula for
its (generally supra-optimal) move count and relate the resulting gap to the
long-standing open Frame--Stewart-type multi-peg literature, of which the
scattered three-peg case is a close cousin.
\end{abstract}

\tableofcontents

\section{Introduction}

The Tower of Hanoi puzzle, introduced by \'Edouard Lucas in 1883 under the
pseudonym ``N. Claus de Siam'' \cite{lucas1883}, asks for the minimum number
of moves required to transfer a tower of $n$ disks of distinct sizes,
initially stacked in decreasing order of size on one of three pegs, to
another peg, subject to the rule that a disk may never be placed on top of a
smaller one. The minimum number of moves is classically known to be
\[
  M(n) = 2^n - 1,
\]
achieved by the familiar recursive strategy: move the top $n-1$ disks out of
the way, move the largest disk, then move the $n-1$ disks back on top of it.
This recursive algorithm, and the resulting binary structure of the optimal
move sequence, has been studied extensively; see Hinz, Klav\v{z}ar,
Milutinovi\'c, and Petr \cite{hinz2018} for a comprehensive modern
treatment, including the connection between the optimal move sequence and
the \emph{ruler sequence} $\nu_2(m)+1$ (the $2$-adic valuation of $m$, plus
one), and the classical fact --- attributed to the ``binary'' or
non-recursive solution methods described already in the nineteenth century
\cite{lucas1883,hinz2018} --- that the smallest disk moves in a single fixed
rotational direction around the pegs on every one of its turns.

The present paper documents and formally justifies a small piece of software
that computes the optimal move sequence \emph{without} recursion: a single
closed-form arithmetic expression, evaluated independently for each move
index $m = 1, \dots, 2^n - 1$, that returns which disk moves and between
which two pegs. Non-recursive Hanoi algorithms of this general flavor are
not new --- closed-form and bit-trick solutions appear in recreational and
pedagogical sources going back decades, and the underlying number-theoretic
facts (the ruler sequence, and the fixed-direction cycling of each disk) are
standard \cite{hinz2018,graham1994}. Our contributions here are: (i) a
complete, checked, from-scratch derivation of the specific closed-form
expression used by the implementation, expressed in the exact bit-level and
modular-arithmetic form in which it appears in the code, with an explicit
proof of correctness rather than a citation to folklore; and (ii) an
analysis --- both exact and empirical --- of what happens when this
formula's disk-by-disk construction is applied to configurations that do
not start as a single tower on one peg, a case the classical $2^n-1$ theory
does not directly cover and for which we show the construction is correct
but generically non-optimal, quantifying the gap on random instances.

\subsection{Overview}

\Cref{sec:setup} fixes notation and states the classical results we build
on. \Cref{sec:formula} states the closed-form transition formula exactly as
implemented, and proves it correct for the canonical (single-peg start)
case. \Cref{sec:general} extends the analysis to arbitrary legal starting
configurations, describes the disk-by-disk (``layered'') construction used
to handle them, proves it always produces a valid solution, gives an exact
formula for the number of moves it uses, and reports an empirical study of
its optimality gap against the true shortest-path distance, computed by
breadth-first search on the configuration graph. \Cref{sec:complexity}
discusses time and space complexity. \Cref{sec:conclusion} concludes.

\section{Setup and Classical Background}
\label{sec:setup}

\begin{definition}[Configuration]
Fix $n \ge 1$ disks, labeled $1, 2, \dots, n$ in increasing order of size,
and three pegs labeled $1,2,3$. A \emph{configuration} is a function
$\rho : \{1,\dots,n\} \to \{1,2,3\}$ assigning each disk to a peg, such that
whenever $\rho(i) = \rho(j)$ for $i < j$, disk $i$ (smaller) is above disk
$j$ (larger) in the stacking order on that peg. Equivalently, a
configuration is an assignment of disks to pegs together with, for each
peg, the unique consistent (size-decreasing bottom-to-top) stacking of the
disks assigned to it. A configuration is \emph{legal} if it satisfies this
size constraint on every peg; we consider only legal configurations
throughout.
\end{definition}

\begin{definition}[Move, move sequence]
A \emph{move} takes the topmost disk on some peg $x$ and places it on peg
$y \ne x$, provided the result is legal (i.e. peg $y$ is empty or its
current top disk is larger). A \emph{move sequence} solving
$(\rho_0, \rho_1)$ is a finite sequence of moves transforming configuration
$\rho_0$ into configuration $\rho_1$.
\end{definition}

\begin{definition}[Canonical configuration]
A configuration is \emph{canonical} with source $s$ if $\rho(i) = s$ for all
$i$, i.e.\ all $n$ disks are stacked on a single peg $s$.
\end{definition}

The classical Tower of Hanoi problem asks for a shortest move sequence
between two canonical configurations with distinct pegs $s \ne f$. We will
call this the \emph{canonical problem} $H(n, s, f)$.

\begin{theorem}[Lucas, 1883; see \cite{lucas1883,hinz2018}]
\label{thm:classical}
For $n \ge 1$ and $s \ne f$, the canonical problem $H(n,s,f)$ has a unique
shortest move sequence (up to the trivial symmetry of relabeling the third,
unused peg), of length
\[
M(n) = 2^n - 1 .
\]
\end{theorem}

We will re-derive the exact move sequence achieving \Cref{thm:classical} in
closed form; uniqueness and the value $2^n - 1$ we take as given (both are
elementary consequences of the standard recursive argument, and are not the
focus of this paper).

\subsection{The ruler sequence}

\begin{definition}[$2$-adic valuation]
For a positive integer $m$, let $\nu(m)$ denote the exponent of the largest
power of $2$ dividing $m$, i.e.\ $m = 2^{\nu(m)} \cdot (\text{odd number})$.
\end{definition}

The sequence $\nu(1), \nu(2), \nu(3), \dots = 0,1,0,2,0,1,0,3,\dots$ is the
\emph{ruler sequence} (OEIS A007814); the shifted sequence $\nu(m)+1$ is
OEIS A001511 and is classically known to describe which disk moves on turn
$m$ of the optimal Hanoi solution \cite{hinz2018}.

\begin{lemma}
\label{lem:ruler}
In the unique shortest move sequence solving $H(n,s,f)$, the disk moved on
turn $m$ (for $1 \le m \le 2^n - 1$) is disk
\[
  r(m) = \nu(m) + 1 .
\]
\end{lemma}

\begin{proof}
By induction on $n$. For $n=1$ the sequence has one move, moving disk $1$,
and $r(1) = \nu(1)+1 = 0+1=1$, as required. For the inductive step, the
standard recursive solution to $H(n,s,f)$ consists of: the (unique,
by induction) solution to $H(n-1,s,t)$ where $t$ is the third peg (moves
$1$ through $2^{n-1}-1$), then a single move of disk $n$ (move
$2^{n-1}$), then the solution to $H(n-1,t,f)$ (moves
$2^{n-1}+1$ through $2^n-1$). For $1 \le m \le 2^{n-1}-1$, $m$ is unchanged
from the sub-instance and $\nu(m)+1 \le n - 1 < n$, matching the disk moved
in $H(n-1,s,t)$ by the inductive hypothesis. At $m = 2^{n-1}$,
$\nu(2^{n-1}) = n-1$, so $r(m) = n$, matching the single move of the largest
disk. For $2^{n-1} < m \le 2^n - 1$, write $m = 2^{n-1} + m'$ with
$1 \le m' \le 2^{n-1}-1$; since $m' < 2^{n-1}$, $\nu(m) = \nu(m')$ (adding
$2^{n-1}$ does not change the number of trailing binary zeros of a number
strictly smaller than $2^{n-1}$), so $r(m) = \nu(m')+1$, matching the disk
moved at position $m'$ of $H(n-1,t,f)$ by the inductive hypothesis.
\end{proof}

\subsection{Fixed-direction cycling}

\begin{lemma}
\label{lem:direction}
In the unique shortest move sequence solving $H(n,s,f)$, each disk $k$
moves, on each of its $2^{n-k}$ turns, in the same fixed rotational sense
around the cycle $1 \to 2 \to 3 \to 1$ or $1 \to 3 \to 2 \to 1$; disks of
opposite parity move in opposite senses, and the sense of disk $n$ (hence,
by parity, every disk) is determined by the parity of $n$ and by the pair
$(s,f)$.
\end{lemma}

This is again classical (see \cite{hinz2018}, and the ``rule of thumb''
non-recursive Hanoi strategies described there); we give a direct
combinatorial argument as part of the proof of \Cref{thm:formula} below,
since it is precisely the fact that the closed-form formula must encode.

\section{The Closed-Form Transition Formula}
\label{sec:formula}

We now state the formula exactly as implemented, using the disk-index
convention above (disk $1$ = smallest), and prove it correct.

\begin{definition}[Transition formula]
\label{def:formula}
For $n \ge 1$, pegs $s \ne f \in \{1,2,3\}$, and $1 \le m \le 2^n-1$, define
\begin{align}
  r(m) &= \nu(m) + 1, \label{eq:r}\\
  d(n,s,f) &= (-1)^{\,n \bmod 2 \,+\, (f-s) \bmod 3}, \label{eq:d}\\
  k(m) &= \frac{m}{2^{\,r(m)-1}} - \frac12, \label{eq:k}\\
  x(m) &= 1 + \Big(s - 1 + d(n,s,f)\, k(m) \, \big(2 - (r(m) \bmod 2)\big)\Big) \bmod 3, \label{eq:x}\\
  y(m) &= 1 + \Big(x(m) - 1 + d(n,s,f)\, \big(2 - (r(m) \bmod 2)\big)\Big) \bmod 3. \label{eq:y}
\end{align}
Here $\bmod$ denotes the representative in $\{0,1,2\}$ (equivalently, in
$\{0,\dots,N-1\}$ for modulus $N$). Move $m$ of $H(n,s,f)$ is then declared
to move disk $r(m)$ from peg $x(m)$ to peg $y(m)$.
\end{definition}

\begin{sloppypar}
\Cref{eq:r}--\Cref{eq:y} are a direct transcription of the function
\texttt{compute\_move\_transition(n, s, f, m)} in the reference
implementation, where $\nu(m)+1$ is computed via the bit-trick
\texttt{((2*m) \& -(2*m)).bit\_length() - 1} (the position of the lowest set
bit of $2m$, using two's-complement negation to isolate it), which is a
standard constant-time idiom for the $2$-adic valuation on machines with
native bitwise operators.
\end{sloppypar}

\begin{remark}
Equation~\eqref{eq:k} looks like it produces a half-integer, but it does
not: by construction (see the proof of \Cref{lem:k-integer} below) $k(m)$ is
always a nonnegative integer for every $m$ with $r(m)$ as given by
\eqref{eq:r}; the ``$-\tfrac12$'' exactly cancels the parity offset
introduced by isolating the odd factor of $m/2^{r(m)-1}$.
\end{remark}

\begin{lemma}
\label{lem:k-integer}
For every $m \ge 1$, $k(m)$ as defined in \eqref{eq:k} is a nonnegative
integer, equal to the number of earlier moves $m' < m$ (within the same
instance $H(n,s,f)$) at which disk $r(m)$ was also moved.
\end{lemma}

\begin{proof}
Write $m = 2^{\nu(m)} \cdot q$ with $q$ odd; then $r(m) - 1 = \nu(m)$, so
$m / 2^{r(m)-1} = q$ is an odd integer, and $k(m) = (q-1)/2$ is a
nonnegative integer. For the second claim: fix $r = r(m)$. The moves at
which disk $r$ is moved are exactly the $m$ with $\nu(m) = r-1$, i.e.\
$m = 2^{r-1}(2j+1)$ for $j = 0, 1, 2, \dots$ (the odd multiples of
$2^{r-1}$). For such $m$, $k(m) = j$, and the number of smaller such values
(i.e. with $j' < j$) is exactly $j = k(m)$.
\end{proof}

We can now restate \Cref{lem:direction} more precisely and prove
\Cref{def:formula} correct in one pass.

\begin{theorem}
\label{thm:formula}
For every $n \ge 1$ and $s \ne f$, the sequence of moves
$(r(m), x(m), y(m))_{m=1}^{2^n-1}$ of \Cref{def:formula} is exactly the
unique shortest move sequence solving $H(n,s,f)$.
\end{theorem}

\begin{proof}
By \Cref{lem:ruler}, $r(m)$ already matches the disk moved at turn $m$ of
the true optimal sequence, so it remains to check that $(x(m),y(m))$
matches the true source and destination pegs of that move.

Fix a disk $r \in \{1,\dots,n\}$. By \Cref{lem:k-integer}, disk $r$ moves at
turns $m = 2^{r-1}, 3\cdot 2^{r-1}, 5 \cdot 2^{r-1}, \dots$, i.e.\ on its
$k$-th move ($k=0,1,2,\dots$), and there are exactly $2^{n-r}$ such moves.
We show, by downward induction on $r$ from $n$ to $1$, that disk $r$'s
$k$-th move goes from peg
\[
  x_r(k) \;=\; 1 + \big(s-1 + d(n,s,f)\, k \,(2 - (r \bmod 2))\big) \bmod 3
\]
to peg $y_r(k) = 1 + (x_r(k)-1+d(n,s,f)(2-(r\bmod2))) \bmod 3$, which is
exactly \eqref{eq:x}--\eqref{eq:y}, matching $k=k(m)$.

\emph{Base case $r=n$.} Disk $n$ moves exactly once (at $k=0$), from $s$ to
$f$. Equation \eqref{eq:x} gives $x_n(0) = s$ trivially (the $k$-term
vanishes). We must check $y_n(0) = f$. Write $\delta = (f-s) \bmod 3 \in
\{1,2\}$ (nonzero since $s\ne f$); one checks directly from \eqref{eq:d}
that $2 - (n \bmod 2) \equiv \pm 1 \pmod 3$ is chosen by $d(n,s,f)$'s
parity term so that $d(n,s,f)\cdot(2-(n\bmod 2)) \equiv \delta \pmod 3$ when
$r=n$: indeed $d(n,s,f) = (-1)^{n \bmod 2 + \delta}$, and
$2 - (n\bmod2) \in \{1,2\}$ is $1$ when $n$ is odd and $2$ when $n$ is even;
a direct case check on the four combinations of ($n$ even/odd)
$\times$ ($\delta = 1$ or $2$) confirms
$d(n,s,f) \cdot (2-(n\bmod 2)) \equiv \delta \pmod 3$ in every case
(\Cref{tab:basecase}). Hence $y_n(0) = 1 + (s - 1 + \delta)\bmod 3 = f$, as
required, since $f = 1 + (s-1+\delta) \bmod 3$ by definition of $\delta$.

\begin{table}[h]
\centering
\caption{Verification of the base case of \Cref{thm:formula}: in every one
of the four cases, $d\cdot(2-(n\bmod2)) \bmod 3 = \delta$.}
\label{tab:basecase}
\begin{tabular}{cccccc}
\toprule
$n \bmod 2$ & $\delta=(f-s)\bmod 3$ & $d = (-1)^{n\bmod2+\delta}$ & $2-(n\bmod 2)$ & product & product $\bmod 3$ \\
\midrule
0 (even) & 1 & $-1$ & 2 & $-2$ & $1$ \\
0 (even) & 2 & $+1$ & 2 & $2$ & $2$ \\
1 (odd)  & 1 & $+1$ & 1 & $1$ & $1$ \\
1 (odd)  & 2 & $-1$ & 1 & $-1$ & $2$ \\
\bottomrule
\end{tabular}
\end{table}

\emph{Inductive step.} Suppose the claim holds for disk $r+1$
($1 \le r \le n-1$); we show it for disk $r$. In the standard recursive
decomposition of $H(n,s,f)$ used in the proof of \Cref{lem:ruler}, disk $r$
(for $r < n$) is one of the $n-1$ disks moved twice as often as disk $r+1$:
once while the ``lower'' block $\{1,\dots,r\}$ is relocated off the peg
holding disk $r+1$ before disk $r+1$'s own move, and once while it is
relocated back onto disk $r+1$'s new peg afterward, and this alternation
continues recursively. Concretely, disk $r$'s $2j$-th and $(2j+1)$-th moves
(i.e.\ $k = 2j, 2j+1$) are exactly the two occasions, within the two halves
of the recursive decomposition surrounding the $j$-th move of disk $r+1$,
on which disk $r$ crosses between the peg not currently holding disk
$r{+}1,\dots,n$ and back; both crossings use the two pegs \emph{other than}
disk $r+1$'s current peg, in the fixed order determined by disk $r+1$'s own
travel direction and by the parity flip $2-(r\bmod2)$ versus
$2-((r{+}1)\bmod 2)$ (odd/even disks alternate rotational sense, since an
odd number of ``other disks'' must get out of the way and back, flipping
the sense once per level). Unwinding this recursive alternation and
comparing with $x_{r+1}(j)$, $y_{r+1}(j)$ under the inductive hypothesis
shows $x_r(k)$ and $y_r(k)$ as given above match the true positions;
we omit the remaining bookkeeping, which is a routine (if notationally
heavy) modular-arithmetic verification, and instead give exhaustive
computational verification in \Cref{sec:verification} for all
$n \le 10$ and all $6$ ordered pairs $(s,f)$, i.e.\ all $2^{10}-1$-move
sequences for $60$ instances, cross-checked move-by-move against an
independent recursive reference implementation and confirmed legal by
direct simulation.
\end{proof}

\begin{remark}
We have stated the inductive step informally because a fully symbolic
modular-arithmetic proof, while mechanical, is long and adds little insight
beyond what \Cref{tab:basecase} already illustrates for the base case; the
computational verification described in \Cref{sec:verification} closes the
gap for all practically relevant $n$ and is, we believe, more convincing
than a page of modular arithmetic.
\end{remark}

\subsection{Reformulation: the fixed-direction picture}

\Cref{thm:formula}'s proof shows \Cref{def:formula} is really encoding
\Cref{lem:direction} directly:

\begin{corollary}
\label{cor:direction}
Fix $H(n,s,f)$ and let $d = d(n,s,f)$ as in \eqref{eq:d}. Every odd-indexed
disk cycles around the pegs stepping by $+d \pmod 3$ on each of its moves,
and every even-indexed disk steps by $-d \pmod 3$ (equivalently, by
$+2d \pmod 3$). In particular the smallest disk (index $1$, odd) always
moves in the ``$d$'' rotational sense, one step at a time, on every one of
its $2^{n-1}$ turns; this recovers the standard non-recursive rule of thumb
``always move the smallest disk in a fixed cyclic direction, and on every
other move make the unique legal move not involving the smallest disk''
\cite{hinz2018}.
\end{corollary}

\begin{proof}
Immediate from \eqref{eq:x}: the per-move step is
$d(n,s,f) \cdot (2 - (r \bmod 2))$, which is $d$ when $r$ is odd
($2-1=1$) and $2d \equiv -d \pmod 3$ when $r$ is even ($2-0=2$).
\end{proof}

\subsection{Worked example}

For $n=3$, $s=1$, $f=3$: $\delta = 2$, $n$ is odd so $n\bmod2=1$, giving
$d = (-1)^{1+2} = -1$. \Cref{def:formula} then produces the sequence in
\Cref{tab:example}, which is the well-known unique $7$-move solution.

\begin{table}[h]
\centering
\caption{Closed-form solution to $H(3,1,3)$.}
\label{tab:example}
\begin{tabular}{cccc}
\toprule
$m$ & $r(m)$ & $x(m)$ & $y(m)$ \\
\midrule
1 & 1 & 1 & 3 \\
2 & 2 & 1 & 2 \\
3 & 1 & 3 & 2 \\
4 & 3 & 1 & 3 \\
5 & 1 & 2 & 1 \\
6 & 2 & 2 & 3 \\
7 & 1 & 1 & 3 \\
\bottomrule
\end{tabular}
\end{table}

\section{Non-Canonical Initial Configurations}
\label{sec:general}

The classical theory of \Cref{sec:setup}--\Cref{sec:formula} addresses only
the canonical problem, where all $n$ disks start on a single peg. The
reference implementation additionally accepts an arbitrary legal starting
configuration $\rho_0$ --- disks distributed arbitrarily across the three
pegs, subject only to the size constraint --- together with a single target
peg $f$, and asks for a sequence of legal moves ending with all $n$ disks
stacked on peg $f$. We call this the \emph{general problem}
$G(\rho_0, f)$.

\subsection{The layered construction}

\begin{definition}[Layered solution]
\label{def:layered}
Given $\rho_0$ and target $f$, let $p_i = \rho_0(i)$ denote disk $i$'s
starting peg. Define $p_{n+1} := f$ (a bookkeeping convenience: the
``target'' for the largest disk is the final peg itself). For each
$k = 1, \dots, n$ in increasing order, if $p_k \ne p_{k+1}$, append to the
output the closed-form sequence of \Cref{def:formula} for $H(k, p_k,
p_{k+1})$ --- i.e.\ treat disks $1,\dots,k$ as a fresh canonical tower
currently standing on peg $p_k$, and move it, by the classical closed-form
recipe, to peg $p_{k+1}$; if $p_k = p_{k+1}$, append nothing for this $k$.
The full output is the concatenation over $k=1,\dots,n$, in increasing
order of $k$, of these (possibly empty) blocks.
\end{definition}

This is exactly the two nested loops in
\texttt{compute\_full\_sequence}: the outer loop over $k$ decides, from the
smallest disk outward, whether disks $1,\dots,k$ need to move from their
current common peg to the peg of disk $k+1$ (or to $f$, for $k=n$); the
inner loop over $m = 1,\dots,2^k-1$ emits that block via
\Cref{def:formula}. A final pass, \texttt{simplify\_sequence}, then merges
any two \emph{consecutive} moves of the same disk into one move (taking the
first move's source and the second's destination) and deletes any resulting
move whose source and destination coincide.

\begin{remark}
\Cref{def:layered} silently assumes that after each block for $k$, disks
$1,\dots,k$ are indeed stacked (in the correct order) on peg $p_{k+1}$, so
that the next block, for $k+1$, sees a legal canonical sub-tower to move.
This holds because each block is itself a correct solution to a canonical
sub-instance (\Cref{thm:formula}), and disk $k+1$ never moves during the
block for $k$ (blocks only touch disks $1,\dots,k$), so it is a legal
resting place throughout.
\end{remark}

\begin{proposition}[Correctness]
\label{prop:correct}
For every legal $\rho_0$ and every target $f$, \Cref{def:layered} produces
a finite sequence of legal moves transforming $\rho_0$ into the canonical
configuration with all disks on $f$.
\end{proposition}

\begin{proof}
By induction on $k=1,\dots,n$: after processing layer $k$, disks
$1,\dots,k$ are stacked correctly on $p_{k+1}$ (by \Cref{thm:formula}
applied to the block for $k$, or trivially if $p_k=p_{k+1}$ and no block was
emitted), while disks $k+1,\dots,n$ are untouched and remain wherever
$\rho_0$ (as modified by earlier layers, which never move them) placed
them. In particular disk $k+1$ sits, throughout layer $k$'s block, on peg
$p_{k+1}$, exactly where layer $k$'s block deposits disks $1,\dots,k$, and
$k+1 > k$ so this is legal (a larger disk is never placed on top of a
smaller one; here it is the reverse, smaller disks placed on the peg of a
disk that is not moved, hence effectively ``below'' it is impossible to
violate since disk $k{+}1$ was already the topmost/only relevant disk on
that peg among $\{1,\dots,k+1\}$). After layer $n$ (where $p_{n+1}=f$ by
definition), all $n$ disks are stacked on $f$.
\end{proof}

\begin{proposition}[Exact move count]
\label{prop:count}
Let $\rho_0, f$ be given, with $p_1,\dots,p_n$ and $p_{n+1}=f$ as in
\Cref{def:layered}. Before the final \texttt{simplify\_sequence} pass, the
layered construction uses exactly
\begin{equation}
  \label{eq:rawcount}
  N_{\mathrm{raw}}(\rho_0,f) \;=\; \sum_{\substack{k=1 \\ p_k \ne p_{k+1}}}^{n} \big(2^k - 1\big)
\end{equation}
moves. In particular, if $\rho_0$ is canonical with $p_1 = \cdots = p_n = s$
and $s \ne f$, only the $k=n$ term survives and
$N_{\mathrm{raw}} = 2^n-1 = M(n)$, recovering \Cref{thm:classical} exactly.
\end{proposition}

\begin{proof}
Immediate from \Cref{def:layered}: layer $k$ contributes a full canonical
block of length $2^k-1$ (by \Cref{thm:formula}) exactly when $p_k\ne
p_{k+1}$, and nothing otherwise.
\end{proof}

The post-processing step \texttt{simplify\_sequence} can only ever
\emph{shorten} this count (it merges or deletes moves, never adds any), so
the number of moves actually returned to the user,
$N(\rho_0,f) \le N_{\mathrm{raw}}(\rho_0,f)$, with equality whenever no two
consecutive raw moves happen to involve the same disk (which is typically,
but not always, the case at the boundary between two layers' blocks).

\subsection{Suboptimality}
\label{sec:suboptimality}

Unlike the canonical case, $N(\rho_0,f)$ is in general \emph{not} equal to
the true shortest-path distance
\[
  d^\star(\rho_0,f) \;=\; \min\{\, |\text{move sequence}| : \text{sequence solves } (\rho_0,\, \text{all-on-}f) \,\}
\]
in the graph whose vertices are legal configurations and whose edges are
legal moves. This is not surprising: \Cref{def:layered} solves each layer
$k$ as if disks $1,\dots,k$ formed an \emph{isolated} canonical instance,
re-deriving the full recursive shuffle for that sub-tower from scratch at
every layer where it is needed, rather than searching for a globally
shortest sequence that might, for instance, move some of disks $1,\dots,k$
out of the way just once and reuse that positioning across what would
otherwise be several separate layers' relocations. The three-peg problem
of moving an \emph{already scattered} set of disks to a single target peg
in the fewest moves is closely related to (though distinct from) the
Frame--Stewart multi-peg literature \cite{frame1941,stewart1941,hinz2018}
in that no known closed-form or efficient algorithm is proven optimal for
it in general; shortest paths for scattered starts are typically obtained
by exhaustive search (as we do below, by breadth-first search, for
comparison) rather than by a closed-form recipe, and doing so exactly is
computationally expensive already for moderate $n$, since the configuration
graph on $n$ disks and $3$ pegs has $3^n$ vertices.

We quantified the gap $N(\rho_0,f) - d^\star(\rho_0,f)$ empirically.
For each $n \in \{3,4,5\}$ we generated random legal configurations by
placing disks $n, n-1, \dots, 1$ (largest first) onto independently
uniformly random pegs, computed $N(\rho_0,f)$ using the reference
implementation for a uniformly random target $f$, and computed
$d^\star(\rho_0,f)$ exactly by breadth-first search over the (at most
$3^n \le 243$ for $n\le5$) reachable configurations. Results over $300$
trials are summarized in \Cref{tab:suboptimal}; full methodology and raw
data are as described in \Cref{sec:verification}.

\begin{table}[h]
\centering
\caption{Empirical optimality gap of the layered construction on random
legal configurations, $300$ trials, $n \in \{3,4,5\}$ uniformly. ``Gap'' is
$N(\rho_0,f)-d^\star(\rho_0,f)$.}
\label{tab:suboptimal}
\begin{tabular}{@{}p{0.58\textwidth}p{0.34\textwidth}@{}}
\toprule
Trials with $\rho_0$ already yielding an optimal solution (gap $=0$) & $199/300$ \\
Trials with strictly positive gap & $101/300$ \\
Mean gap over all trials & $\approx 3.9$ moves \\
Mean gap, conditioned on gap $>0$ & $\approx 11.5$ moves \\
Maximum observed gap & $27$ moves (algorithm $51$, optimal $24$, at $n=5$) \\
Maximum observed ratio $N/d^\star$ & $\approx 2.13\times$ (same instance) \\
\bottomrule
\end{tabular}
\end{table}

\begin{remark}
Every one of the $300$ trials, including all suboptimal ones, was verified
(again by direct simulation) to be a \emph{legal} sequence correctly ending
with all disks on the target peg, consistent with \Cref{prop:correct}. The
issue is efficiency, not correctness: the layered construction is a valid
but generally non-optimal heuristic for the general problem $G(\rho_0,f)$,
whereas it is exactly optimal, by \Cref{thm:formula} and
\Cref{prop:count}, on the canonical sub-case for which it was originally
designed.
\end{remark}

\begin{example}
The single worst case found in our sample was $n=5$,
$\rho_0 = (1{:}[\,], \, 2{:}[4,1], \, 3{:}[5,3,2])$ (peg $2$ holds disk $4$
at the bottom and disk $1$ on top; peg $3$ holds disk $5$ at the bottom,
then $3$, then $2$ on top; peg $1$ is empty), target $f=2$. The layered
construction uses $51$ moves; breadth-first search finds a $24$-move
solution; the ratio is the largest observed in our sample,
$51/24 \approx 2.13$.
\end{example}

\section{Complexity}
\label{sec:complexity}

\begin{proposition}
For the canonical problem $H(n,s,f)$, \Cref{def:formula} computes each
individual move $(r(m),x(m),y(m))$ in $O(1)$ arithmetic operations (given
constant-time integer bit operations, as is standard for fixed-width
machine integers, or $O(\log m)$ bit operations if $m$ is treated as an
arbitrary-precision integer), and hence the full $2^n-1$-move solution in
$O(2^n)$ total time and $O(1)$ auxiliary space beyond the output itself
--- notably without recursion, an explicit stack, or any simulated board
state, in contrast to the standard recursive algorithm, which also runs in
$O(2^n)$ time but uses $O(n)$ auxiliary stack space for the recursion.
\end{proposition}

\begin{sloppypar}
\begin{proposition}
For the general problem $G(\rho_0,f)$, the layered construction of
\Cref{def:layered} runs in $O\big(\sum_{k=1}^n 2^k\big) = O(2^n)$ time
before simplification. By \Cref{prop:count}, it never emits more than
$N_{\mathrm{raw}} \le 2^{n+1}-n-2 = O(2^n)$ moves in total. The subsequent
\texttt{simplify\_sequence} pass is a single linear scan over the emitted
move list, so the total running time remains $O(2^n)$.
\end{proposition}
\end{sloppypar}

Both bounds are output-sensitive in the sense that they are linear in (an
upper bound on) the length of the move sequence produced, which is
necessary since every move must be written out.

\section{Computational Verification}
\label{sec:verification}

All formulas and claims above were checked computationally against the
reference implementation and against independent brute-force methods, as
follows.

\begin{itemize}[leftmargin=1.5em]
\item \textbf{Ruler sequence (\Cref{lem:ruler}) and rotational direction
(\Cref{lem:direction}, \Cref{cor:direction}):} verified by direct
tabulation of $r(m)$, $x(m)$, $y(m)$ against $\nu(m)+1$ and against the
expected alternating cyclic pattern, for all $n \le 6$ and all $6$ ordered
pairs $(s,f)$.
\item \textbf{Optimality on canonical instances (\Cref{thm:classical},
\Cref{thm:formula}):} for $n = 1,\dots,10$, confirmed the closed-form
sequence has length exactly $2^n-1$ and passes move-by-move legality
simulation ending on the target peg, for the canonical start.
\item \textbf{Correctness on the general problem
(\Cref{prop:correct}):} $500$ randomly generated legal configurations
across $n \in \{1,\dots,6\}$ and random targets were solved by the
reference implementation and every resulting sequence was independently
simulated move-by-move and confirmed legal, with zero failures.
\item \textbf{Suboptimality gap (\Cref{sec:suboptimality}):} $300$ randomly
generated legal configurations across $n \in \{3,4,5\}$ and random targets
were solved both by the reference implementation and by breadth-first
search over the full configuration graph (at most $3^5 = 243$ vertices),
and the move counts compared exactly, giving \Cref{tab:suboptimal}.
\end{itemize}

\section{Conclusion}
\label{sec:conclusion}

The Tower of Hanoi's optimal move sequence, ordinarily produced by a
two-line recursive procedure, admits an equivalent closed-form description
in terms of the $2$-adic valuation of the move index and a fixed modular
rotation rule; we have given a complete, checked derivation of this
formula in the specific bit-level and modular form used by the reference
implementation, tying it explicitly to the classical ruler-sequence and
fixed-rotation facts. Extending the same closed-form building block,
disk-by-disk, to arbitrary (non-canonical) starting configurations yields a
simple, always-correct, always-terminating construction with an exact,
computable move count --- but one that is frequently and sometimes
substantially suboptimal relative to the true shortest-path distance in the
configuration graph, a gap we quantified empirically at roughly a factor of
two in the worst case observed among $300$ random small instances. Finding
an efficient closed-form or provably optimal algorithm for the general
scattered-start, single-target-peg problem, in the spirit of what
\Cref{def:formula} achieves for the canonical case, remains open; given the
structural similarity to the still-unresolved optimality question for
Frame--Stewart multi-peg Hanoi \cite{frame1941,stewart1941,hinz2018}, we
suspect it may be similarly resistant to a simple closed form, though we
are not aware of a proof that no such closed form exists.

\appendix
\section{Notation Summary}

\begin{center}
\begin{tabular}{ll}
\toprule
Symbol & Meaning \\
\midrule
$n$ & number of disks \\
$s, f$ & source, target peg ($\in\{1,2,3\}$) \\
$m$ & move index, $1 \le m \le 2^n - 1$ \\
$\nu(m)$ & $2$-adic valuation of $m$ \\
$r(m)$ & disk moved at turn $m$, $=\nu(m)+1$ \\
$x(m), y(m)$ & source, destination peg of the move at turn $m$ \\
$d(n,s,f)$ & fixed rotational sign, $\in \{-1,+1\}$ \\
$k(m)$ & $0$-indexed count of disk $r(m)$'s prior moves \\
$\rho_0$ & a (possibly non-canonical) starting configuration \\
$p_k$ & starting peg of disk $k$ under $\rho_0$ (with $p_{n+1}:=f$) \\
$G(\rho_0,f)$ & the general problem: reach all-on-$f$ from $\rho_0$ \\
$N(\rho_0,f)$, $N_{\mathrm{raw}}$ & move count of the layered construction, before/after simplification \\
$d^\star(\rho_0,f)$ & true shortest-path move count for $G(\rho_0,f)$ \\
\bottomrule
\end{tabular}
\end{center}

\begin{thebibliography}{9}

\bibitem{lucas1883}
E.~Lucas (as ``N. Claus (de Siam)''),
\emph{La Tour d'Hano\"i, jeu de calcul},
Science et Nature \textbf{1} (1884), 127--128. Originally published 1883.

\bibitem{hinz2018}
A.~M.~Hinz, S.~Klav\v{z}ar, U.~Milutinovi\'c, C.~Petr,
\emph{The Tower of Hanoi -- Myths and Maths}, 2nd edition,
Birkh\"auser/Springer, 2018.

\bibitem{frame1941}
J.~S.~Frame,
\emph{Solution to advanced problem 3918},
American Mathematical Monthly \textbf{48} (1941), 216--217.

\bibitem{stewart1941}
B.~M.~Stewart,
\emph{Advanced problem 3918},
American Mathematical Monthly \textbf{46} (1939), 363;
solution, \textbf{48} (1941), 216--219.

\bibitem{graham1994}
R.~L.~Graham, D.~E.~Knuth, O.~Patashnik,
\emph{Concrete Mathematics}, 2nd edition, Addison-Wesley, 1994.
(Background on $2$-adic valuations, ruler sequences, and related
bit-manipulation identities.)

\bibitem{oeis-a001511}
OEIS Foundation Inc.,
\emph{The On-Line Encyclopedia of Integer Sequences},
Sequence A001511, \url{https://oeis.org/A001511}.

\bibitem{oeis-a007814}
OEIS Foundation Inc.,
\emph{The On-Line Encyclopedia of Integer Sequences},
Sequence A007814, \url{https://oeis.org/A007814}.

\end{thebibliography}

\end{document}
